% Source: Gerzensee "Real Analysis" slides 5-25 (inner product, Riesz
% representation for R^n, norm, Cauchy-Schwarz, distance, the nesting of inner
% product / normed / metric spaces, Banach and Hilbert spaces).
% Supplementary: Fukuda (2026), Sec. 7.1 (pp. 201-214), Def. 7.12-7.13 (p. 293).
% External references actually cited in this chapter: LeRoy-Werner (2000).

\chapter{Inner Products, Norms, and Metric Spaces}\label{ch:metric}

\section{Motivation: Geometry Without Coordinates}\label{sec:met-motivation}

The Gerzensee course opens its real-analysis block with the inner product and
immediately lists what it buys: length, angle, and orthogonality in $\R^n$ and in
function spaces; the projection theorem, which delivers both ordinary least
squares and conditional expectation; and the separation theorems, which deliver
the second welfare theorem and the existence of state prices in a no-arbitrage
economy \autocite[slide 3]{fukuda2026analysis}. All of these are geometric
statements, and all of them survive the passage from $\R^n$ to spaces of
functions and of random variables --- provided the geometry is set up
abstractly.

The structures form a strict hierarchy:
\[
	\text{inner product spaces}\ \subsetneq\
	\text{normed spaces}\ \subsetneq\
	\text{metric spaces},
\]
reproduced on slide~25 of \autocite{fukuda2026analysis}. Each step forgets
something. A metric gives distance but no notion of adding points; a norm gives
length and a vector structure but no notion of angle; an inner product gives
angle as well. Knowing which structure a theorem needs is what tells us how far
it generalises, and this chapter states each result at the weakest level at which
it is true.

\section{Inner Product Spaces}\label{sec:met-ip}

\begin{definition}[Inner product]\label{def:inner-product}
	Let $V$ be a real vector space. An \dfn{inner product} on $V$ is a map
	$\ip{\cdot}{\cdot}\colon V\times V\to\R$ such that for all $x,y,z\in V$ and
	$\lambda\in\R$:
	\begin{enumerate}[(i)]
		\item $\ip{x}{x}\geq0$, with $\ip{x}{x}=0$ if and only if $x=0$
		      (positive definiteness);
		\item $\ip{x}{y}=\ip{y}{x}$ (symmetry);
		\item $\ip{x}{\lambda y}=\lambda\ip{x}{y}$ and
		      $\ip{x}{y+z}=\ip{x}{y}+\ip{x}{z}$ (linearity in the second argument;
		      with (ii) this gives linearity in each argument separately).
	\end{enumerate}
	The pair $(V,\ip{\cdot}{\cdot})$ is an \dfn{inner product space}.
	\nidx{$\langle\cdot,\cdot\rangle$}{$\ip{x}{y}$, inner product}
\end{definition}

These are exactly the five properties verified in
\autocite[slides 5--6]{fukuda2026analysis} for the Euclidean inner product
$\ip{a}{x}=\sum_{i=1}^n a_ix_i$ on $\R^n$; the abstraction consists in taking
them as the definition.

\begin{eg}[Inner product spaces used in this volume]\label{eg:ip-spaces}
	\begin{enumerate}[(i)]
		\item $\R^n$ with $\ip{x}{y}=x^{\transp}y=\sum_ix_iy_i$.
		\item $\R^n$ with $\ip{x}{y}_W=x^{\transp}Wy$ for a symmetric $W\succ0$. This is
		      an inner product because $W$ is positive definite
		      (\autoref{prop:definiteness-eigen}), and it is the metric under which
		      generalised least squares is an orthogonal projection.
		\item $\Lspace{2}(\R)=\bigl\{f\colon\R\to\R \text{ measurable}:
			      \int_{-\infty}^{\infty}\abs{f(x)}^2\,\diff x<\infty\bigr\}$ with
		      $\ip{f}{g}=\int_{-\infty}^{\infty}f(x)g(x)\,\diff x$.
		\item The space of square-integrable random variables on a probability space
		      $(\Omega,\mathcal{F},\Prob)$, with $\ip{X}{Y}=\E[XY]$; or, on the subspace
		      of centred variables, $\ip{X}{Y}=\E[(X-\E X)(Y-\E Y)]=\Cov(X,Y)$.
	\end{enumerate}
	Items (iii) and (iv) are stated in \autocite[slide 7]{fukuda2026analysis}.
\end{eg}

\begin{remark}[The qualification the slides leave out]\label{rem:l2-quotient}
	On $\Lspace{2}$ the map $\ip{f}{g}=\int fg$ satisfies (ii) and (iii) but only
	half of (i): $\ip{f}{f}=0$ forces $f=0$ almost everywhere, not everywhere. The
	slide's hedge that ``roughly, the inner product satisfies the properties''
	\autocite[slide 7]{fukuda2026analysis} refers to this. The remedy is
	standard: work with the quotient of $\Lspace{2}$ by the subspace of functions
	vanishing almost everywhere, writing $L^2$ for the quotient, whose elements are
	equivalence classes under the relation of \autoref{prop:equivalence-partition}.
	On $L^2$, positive definiteness holds. In the random-variable case (iv) the same
	convention identifies random variables that agree $\Prob$-almost surely, and it
	is what makes conditional expectation in \autoref{ch:probability} a
	well-defined element rather than a function defined up to null sets.
\end{remark}

\begin{lemma}[Cauchy--Schwarz inequality]\label{lem:cauchy-schwarz}
	Let $(V,\ip{\cdot}{\cdot})$ be an inner product space and set
	$\norm{x}=\sqrt{\ip{x}{x}}$. Then
	\[
		\abs{\ip{x}{y}}\leq\norm{x}\,\norm{y}\qquad\text{for all }x,y\in V,
	\]
	with equality if and only if $x$ and $y$ are linearly dependent.
\end{lemma}

\begin{proof}
	If $y=0$ both sides vanish and $x,y$ are dependent. Assume $y\neq0$ and set
	$\lambda\eqdef\ip{x}{y}/\ip{y}{y}$. Then
	\[
		0\leq\norm{x-\lambda y}^2=\ip{x-\lambda y}{x-\lambda y}
		=\norm{x}^2-2\lambda\ip{x}{y}+\lambda^2\norm{y}^2
		=\norm{x}^2-\frac{\ip{x}{y}^2}{\norm{y}^2},
	\]
	which rearranges to $\ip{x}{y}^2\leq\norm{x}^2\norm{y}^2$.

	Equality holds if and only if $\norm{x-\lambda y}=0$, that is
	$x=\lambda y=\frac{\ip{x}{y}}{\ip{y}{y}}y$, which is linear dependence; and
	conversely if $x=\mu y$ then both sides equal $\abs{\mu}\norm{y}^2$.
\end{proof}

The choice of $\lambda$ is the one that makes $x-\lambda y$ orthogonal to $y$, so
the proof is the projection argument of \autoref{ch:projection} in miniature.
This is the content of the geometric derivation on
\autocite[slide 17]{fukuda2026analysis}, where the condition
$\ip{\lambda x}{y-\lambda x}=0$ is solved for $\lambda$.

\begin{definition}[Orthogonality and angle]\label{def:orthogonality}
	Vectors $x,y$ in an inner product space are \dfn{orthogonal}, written
	$x\perp y$, if $\ip{x}{y}=0$. For non-zero $x,y$ the \dfn{angle} between them is
	the unique $\theta\in[0,\pi]$ with
	$\cos\theta=\ip{x}{y}/(\norm{x}\norm{y})$, which is well defined precisely
	because \autoref{lem:cauchy-schwarz} places the ratio in $[-1,1]$.
\end{definition}

\begin{eg}[Correlation is a cosine]\label{eg:correlation-cosine}
	Let $x,y\in\R^n$ with sample means $\mu_x,\mu_y$, and put
	$\tilde x=x-\mu_x\mathbf{1}$, $\tilde y=y-\mu_y\mathbf{1}$ where
	$\mathbf 1=(1,\dots,1)^{\transp}$. With
	$\sigma_{x,y}=\frac1n\ip{\tilde x}{\tilde y}$,
	$\sigma_x^2=\frac1n\norm{\tilde x}^2$ and $\sigma^2_y=\frac1n\norm{\tilde y}^2$,
	the Pearson correlation coefficient is
	\[
		\rho=\frac{\sigma_{x,y}}{\sigma_x\sigma_y}
		=\frac{\ip{\tilde x}{\tilde y}}{\norm{\tilde x}\,\norm{\tilde y}}
		=\cos\theta .
	\]
	\autoref{lem:cauchy-schwarz} therefore gives $\abs{\rho}\leq1$ immediately, and
	the equality case gives $\abs{\rho}=1$ if and only if $\tilde y=c\tilde x$ for
	some $c\in\R$, that is
	\[
		y_i-\mu_y=\frac{\sigma_{x,y}}{\sigma_x^2}(x_i-\mu_x)
		\qquad\text{for all }i,
	\]
	since $c$ must equal $\ip{\tilde x}{\tilde y}/\norm{\tilde x}^2$. This settles
	Exercise~7.6 of \autocite[slide 18]{fukuda2026analysis} and identifies the
	constant as the population regression slope of \autoref{ch:projection}. The same
	computation with $\ip{X}{Y}=\Cov(X,Y)$ shows that the correlation of two random
	variables is the cosine of the angle between them in $L^2$.
\end{eg}

\subsection{Riesz representation in \texorpdfstring{$\R^n$}{Rn}}

\begin{theorem}[Riesz representation theorem for $\R^n$]\label{thm:riesz-rn}
	\begin{enumerate}[(i)]
		\item For each $a\in\R^n$ the map $f_a\colon\R^n\to\R$, $f_a(x)=\ip{a}{x}$, is
		      linear.
		\item Conversely, for every linear $f\colon\R^n\to\R$ there is a unique
		      $a\in\R^n$ with $f(x)=\ip{a}{x}$ for all $x\in\R^n$; explicitly
		      $a=(f(e^1),\dots,f(e^n))^{\transp}$.
	\end{enumerate}
\end{theorem}

\begin{proof}
	(i) is linearity of the inner product in its second argument.

	(ii) Existence: writing $x=\sum_ix_ie^i$ and using linearity,
	$f(x)=\sum_ix_if(e^i)=\ip{a}{x}$ with $a$ as displayed. Uniqueness: if
	$\ip{a}{x}=\ip{a'}{x}$ for all $x$, then $\ip{a-a'}{x}=0$ for all $x$; taking
	$x=a-a'$ gives $\norm{a-a'}^2=0$, hence $a=a'$ by positive definiteness.
\end{proof}

This is Theorem~7.1 of \textcite{fukuda2026notes}; the two-dimensional picture is
\autocite[slide 9]{fukuda2026analysis}. It is the special case $m=1$ of
\autoref{thm:matrix-rep}, but stated this way it says something the matrix
version does not: \emph{a linear functional is the same thing as a vector}. Three
economic readings follow.

\begin{eg}[Expected utility as an inner product]\label{eg:eu-riesz}
	Let $\Omega=\{\omega_1,\dots,\omega_n\}$ and let $p\in\R^n$ be a probability
	vector. Expected utility $\E[u]=\sum_ip_iu_i=\ip{u}{p}$ is linear in the
	probability vector $p$ for fixed utility index $u$, and linear in $u$ for fixed
	$p$ \autocite[slide 8]{fukuda2026analysis}. Read in the first way,
	\autoref{thm:riesz-rn} says that \emph{any} functional on lotteries that is
	linear in probabilities has an expected-utility representation, with the
	representing vector $a$ being the utility index. The von Neumann--Morgenstern
	theorem is the statement that the behavioural axioms force linearity; the Riesz
	theorem then supplies the index.
\end{eg}

\begin{eg}[State prices and the stochastic discount factor]\label{eg:state-prices}
	Let there be $n$ states and let an asset be a payoff vector $z\in\R^n$. Suppose
	the pricing functional $\pi\colon\R^n\to\R$ is linear --- which is exactly the
	law of one price plus portfolio additivity. By \autoref{thm:riesz-rn} there is a
	unique $q\in\R^n$ with $\pi(z)=\ip{q}{z}=\sum_iq_iz_i$; the $q_i$ are
	\dfnas{state price}{state prices}. If moreover $\pi(z)>0$ whenever
	$z\geq0$, $z\neq0$ --- absence of arbitrage --- then $q\in\R^n_{++}$, and
	writing $q_i=p_im_i$ for a probability $p$ gives
	$\pi(z)=\E[mz]$, exhibiting $m$ as the \dfn{stochastic discount factor}. The
	step from ``no arbitrage'' to ``a strictly positive linear functional exists''
	is not Riesz but a separation theorem, proved in \autoref{ch:convex}; Riesz
	converts the functional into a vector. See
	\textcite[Ch.~2]{leroy2000principles}.
\end{eg}

\begin{eg}[The total differential]\label{eg:total-differential}
	For $f\colon\R^n\to\R$ differentiable at $x$, the differential
	\[
		\diff y=\sum_{i=1}^{n}\pderiv{f}{x_i}(x)\,\diff x_i
	\]
	is a linear functional of
	the displacement $\diff x$. By \autoref{thm:riesz-rn} it is represented by a unique
	vector, and that vector is by definition the gradient:
	$\diff y=\ip{\nabla f(x)}{\diff x}$ \autocite[slide 10]{fukuda2026analysis}. This is the
	cleanest statement of what a gradient is, and it explains why the gradient is a
	column vector while the derivative $Df(x)$ is a row: the derivative is the
	functional, the gradient is its Riesz representative.
\end{eg}

\section{Normed Spaces}\label{sec:met-norm}

\begin{definition}[Norm]\label{def:norm}
	Let $V$ be a real vector space. A \dfn{norm} is a map
	$\norm{\cdot}\colon V\to\R$ such that for all $x,y\in V$, $\lambda\in\R$:
	\begin{enumerate}[(i)]
		\item $\norm{x}\geq0$;
		\item $\norm{x}=0$ if and only if $x=0$;
		\item $\norm{\lambda x}=\abs{\lambda}\norm{x}$ (absolute homogeneity);
		\item $\norm{x+y}\leq\norm{x}+\norm{y}$ (triangle inequality).
	\end{enumerate}
	The pair $(V,\norm{\cdot})$ is a \dfn{normed space}.
	\nidx{$\|\cdot\|$}{$\norm{x}$, norm}
\end{definition}

\begin{proposition}[An inner product induces a norm]\label{prop:ip-induces-norm}
	If $\ip{\cdot}{\cdot}$ is an inner product on $V$, then
	$\norm{x}\eqdef\sqrt{\ip{x}{x}}$ is a norm.
\end{proposition}

\begin{proof}
	(i) and (ii) are positive definiteness; (iii) follows from bilinearity, since
	$\norm{\lambda x}^2=\lambda^2\norm{x}^2$. For (iv),
	\[
		\norm{x+y}^2=\norm{x}^2+2\ip{x}{y}+\norm{y}^2
		\leq\norm{x}^2+2\norm{x}\norm{y}+\norm{y}^2=(\norm{x}+\norm{y})^2,
	\]
	using \autoref{lem:cauchy-schwarz}; taking square roots gives the claim.
\end{proof}

\begin{proposition}[Reverse triangle inequality]\label{prop:reverse-triangle}
	In any normed space, $\bigl|\,\norm{x}-\norm{x_0}\,\bigr|\leq\norm{x-x_0}$.
\end{proposition}

\begin{proof}
	$\norm{x}=\norm{(x-x_0)+x_0}\leq\norm{x-x_0}+\norm{x_0}$ gives
	$\norm{x}-\norm{x_0}\leq\norm{x-x_0}$. Exchanging $x$ and $x_0$ and using
	$\norm{x_0-x}=\norm{x-x_0}$ gives $\norm{x_0}-\norm{x}\leq\norm{x-x_0}$.
	Together these bound the absolute value.
\end{proof}

This is Exercise~7.4 of \autocite[slide 20]{fukuda2026analysis}. It is used at
once in \autoref{ch:continuity} to show that the norm is a continuous function
--- indeed \autoref{prop:reverse-triangle} says it is Lipschitz with constant
$1$.

\begin{eg}[Norms used in this volume]\label{eg:norms}
	On $\R^n$: the Euclidean norm $\norm{x}_2=\bigl(\sum_ix_i^2\bigr)^{1/2}$,
	induced by the Euclidean inner product; the $\ell^1$ norm
	$\norm{x}_1=\sum_i\abs{x_i}$; the supremum norm
	$\norm{x}_\infty=\max_i\abs{x_i}$. On $\Cspace{[a,b]}{\R}$ and on
	$\Bspace{X}{\R}$: the \dfn{supremum norm}
	\[
		\norm{f}_\infty\eqdef\sup_{x}\abs{f(x)},
	\]
	which is the norm used for value functions in \autoref{ch:dynamic}
	\autocite[slides 21--22, 73]{fukuda2026analysis}.
\end{eg}

\begin{proposition}[The supremum norm is a norm]\label{prop:sup-norm}
	Let $X$ be a non-empty set and $\Bspace{X}{\R}$ the set of bounded functions
	$f\colon X\to\R$. Then $\Bspace{X}{\R}$ is a vector space under pointwise
	operations and $\norm{f}_\infty=\sup_{x\in X}\abs{f(x)}$ is a norm on it.
\end{proposition}

\begin{proof}
	Boundedness is preserved by sums and scalar multiples, so $\Bspace{X}{\R}$ is a
	subspace of $\R^X$, and $\norm{f}_\infty<\infty$ for every $f$ in it. Property
	(i) is clear. For (ii), $\norm{f}_\infty=0$ forces $\abs{f(x)}=0$ for every $x$,
	that is $f=0$ --- note that here, unlike in \autoref{rem:l2-quotient}, no
	almost-everywhere qualification arises. For (iii),
	$\sup_x\abs{\lambda f(x)}=\abs{\lambda}\sup_x\abs{f(x)}$. For (iv), for each $x$,
	$\abs{f(x)+g(x)}\leq\abs{f(x)}+\abs{g(x)}\leq\norm{f}_\infty+\norm{g}_\infty$;
	as the right-hand side does not depend on $x$, it is an upper bound for
	$\{\abs{(f+g)(x)}\}$, so by \autoref{prop:sup-char} it dominates the supremum.
\end{proof}

\begin{remark}[Not every norm comes from an inner product]\label{rem:parallelogram}
	A norm is induced by an inner product if and only if it satisfies the
	\dfn{parallelogram law}
	\[
		\norm{x+y}^2+\norm{x-y}^2=2\norm{x}^2+2\norm{y}^2 .
	\]
	Necessity is immediate by expanding both sides. The supremum norm fails it: in
	$\R^2$ with $x=(1,0)$, $y=(0,1)$ we get
	$\norm{x+y}_\infty=\norm{x-y}_\infty=1$, so the left side is $2$ while the right
	side is $4$. Hence $\Bspace{X}{\R}$ with the supremum norm carries no compatible
	inner product, and no notion of angle. This is why the value-function arguments
	of \autoref{ch:dynamic} use the contraction mapping theorem, which needs only a
	metric, rather than a projection argument, which needs an inner product.
\end{remark}

\section{Metric Spaces}\label{sec:met-metric}

\begin{definition}[Metric]\label{def:metric}
	Let $X$ be a non-empty set. A \dfn{metric} (or distance) on $X$ is a map
	$d\colon X\times X\to\R$ such that for all $x,y,z\in X$:
	\begin{enumerate}[(i)]
		\item $d(x,y)\geq0$;
		\item $d(x,y)=0$ if and only if $x=y$;
		\item $d(x,y)=d(y,x)$ (symmetry);
		\item $d(x,z)\leq d(x,y)+d(y,z)$ (triangle inequality).
	\end{enumerate}
	The pair $(X,d)$ is a \dfn{metric space}. \nidx{$d(\cdot,\cdot)$}{$d(x,y)$, metric}
\end{definition}

\begin{proposition}[A norm induces a metric]\label{prop:norm-induces-metric}
	If $(V,\norm{\cdot})$ is a normed space, then $d(x,y)\eqdef\norm{x-y}$ is a
	metric on $V$.
\end{proposition}

\begin{proof}
	(i)--(iii) are immediate from \autoref{def:norm}(i)--(iii), using
	$\norm{y-x}=\norm{-(x-y)}=\norm{x-y}$. For (iv),
	$d(x,z)=\norm{(x-y)+(y-z)}\leq\norm{x-y}+\norm{y-z}=d(x,y)+d(y,z)$.
\end{proof}

Combining \autoref{prop:ip-induces-norm} and \autoref{prop:norm-induces-metric}
gives the hierarchy of \autocite[slide 25]{fukuda2026analysis}: every inner
product space is a normed space and every normed space is a metric space. Both
inclusions are strict --- \autoref{rem:parallelogram} shows the first, and the
next example the second.

\begin{eg}[A metric that comes from no norm]\label{eg:discrete-metric}
	On any non-empty set $X$ the \dfn{discrete metric}
	\[
		d(x,y)=\begin{cases}0 & x=y\\ 1 & x\neq y\end{cases}
	\]
	is a metric. If $X$ is a vector space with more than one point, $d$ is induced
	by no norm: a norm satisfies $\norm{2x}=2\norm{x}$, so
	$d(2x,0)=2d(x,0)$, but the discrete metric gives $d(2x,0)=d(x,0)=1$ for any
	$x\neq0$. More generally, a metric on a vector space comes from a norm if and
	only if it is translation invariant and absolutely homogeneous.

	The discrete metric is not a curiosity: under it every subset is open, every
	sequence that converges is eventually constant, and the compact sets are exactly
	the finite ones. It is the standard source of counterexamples in
	\autoref{ch:topology}.
\end{eg}

\begin{eg}[Metrics on economic objects]\label{eg:econ-metrics}
	\begin{enumerate}[(i)]
		\item Consumption bundles in $\R^n$: Euclidean metric.
		\item Consumption \emph{plans} $c\colon[0,\infty)\to\R$ or
		      $(c_t)_{t\in\N_0}$: supremum metric on $\Bspace{\cdot}{\R}$. The slide's
		      remark that one chooses ``a consumption function
		      $c\colon[0,\infty)\to\R$ instead of $(c_1,c_2)\in\R^2$''
		      \autocite[slide 4]{fukuda2026analysis} is precisely the passage from (i)
		      to (ii), and it is the passage that makes the infinite-dimensional theory
		      necessary.
		\item Value functions $V\colon\R^n\to\R$: supremum metric, under which the
		      Bellman operator is a contraction (\autoref{thm:blackwell}).
		\item Probability distributions: many metrics are in use (total variation,
		      Kolmogorov--Smirnov, Wasserstein), and they are not equivalent in infinite
		      dimensions; which one is meant must always be stated.
	\end{enumerate}
\end{eg}

\begin{definition}[Open ball, bounded set]\label{def:ball-bounded}
	Let $(X,d)$ be a metric space, $x_0\in X$ and $\lambda>0$. The \dfn{open ball}
	with centre $x_0$ and radius $\lambda$ is
	\[
		\ball{\lambda}{x_0}\eqdef\{x\in X : d(x_0,x)<\lambda\}.
	\]
	A set $A\subseteq X$ is \dfn{bounded} if $A=\emptyset$ or $A\subseteq
		\ball{r}{x_0}$ for some $x_0\in X$ and $r>0$.
	\nidx{$B_\lambda(x_0)$}{$\ball{\lambda}{x_0}$, open ball}
\end{definition}

\begin{proposition}\label{prop:bounded-equiv}
	A set $A\subseteq\R^n$ is bounded in the sense of \autoref{def:ball-bounded} if
	and only if there exists $\alpha>0$ with $\norm{x}\leq\alpha$ for all $x\in A$.
\end{proposition}

\begin{proof}
	If $\norm{x}\leq\alpha$ on $A$, then $A\subseteq\ball{\alpha+1}{0}$. Conversely
	if $A\subseteq\ball{r}{x_0}$ then for $x\in A$,
	$\norm{x}\leq\norm{x-x_0}+\norm{x_0}<r+\norm{x_0}$, so
	$\alpha=r+\norm{x_0}$ works.
\end{proof}

This settles Exercise~7.20 of \autocite[slide 28]{fukuda2026analysis}. The
economic instance given there is the budget set: for $p\in\R^n_{++}$ and $m>0$,
every $x\in\budget(p,m)$ satisfies $x_i\leq m/p_i$, so
$\norm{x}_\infty\leq m/\min_ip_i$ and $\budget(p,m)$ is bounded. If some
$p_i=0$, the bound fails and so does the conclusion --- the substantive reason
for requiring strictly positive prices, flagged already in
\autoref{eg:consumer-setup}.

\section{Complete Normed and Inner Product Spaces}\label{sec:met-banach}

The definitions below are stated here for completeness of the hierarchy; the
notion of a Cauchy sequence on which they rest is developed in
\autoref{ch:topology} and their consequences in \autoref{ch:completeness}.

\begin{definition}[Banach and Hilbert spaces]\label{def:banach-hilbert}
	A normed space that is complete as a metric space
	(\autoref{def:complete-space}) is a \dfn{Banach space}. An inner product space
	that is complete in the induced norm is a \dfn{Hilbert space}.
\end{definition}

These are Definitions~7.12 and~7.13 of \textcite{fukuda2026notes}, given on
slide~77 of \autocite{fukuda2026analysis}. The two examples that matter here are
$L^2$, which is a Hilbert space and is where \autoref{ch:projection} operates,
and $\Bspace{X}{\R}$ with the supremum norm, which is a Banach space but not a
Hilbert space by \autoref{rem:parallelogram}, and is where \autoref{ch:dynamic}
operates.

\begin{remark}[Why completeness is stated at the level of the metric]\label{rem:completeness-level}
	Completeness is a metric notion: it makes sense in $(X,d)$ without any vector
	structure. Banach and Hilbert spaces are therefore not new kinds of object but
	conjunctions --- ``normed and complete'', ``inner product and complete''. This
	matters when one restricts attention to a subset: a closed subset of a complete
	metric space is complete (\autoref{prop:closed-subset-complete}), and a closed
	subset of a Banach space need not be a subspace, but is still a complete metric
	space. The contraction mapping theorem applies to it regardless, which is what
	allows \autoref{ch:dynamic} to work on the set of bounded, continuous,
	\emph{increasing and concave} functions --- a closed subset, not a subspace.
\end{remark}

\section{Chapter Summary}\label{sec:met-summary}

An inner product supplies length, angle and orthogonality; the Cauchy--Schwarz
inequality is the statement that the resulting cosine lies in $[-1,1]$, and its
equality case characterises linear dependence, which in the statistical reading
is perfect correlation. Every linear functional on $\R^n$ is an inner product
with a unique vector (\autoref{thm:riesz-rn}); this identification is what turns
linearity in probabilities into an expected-utility index, linearity of a pricing
functional into a vector of state prices, and the differential into the gradient.
A norm forgets angle, a metric forgets the vector structure as well; both
inclusions in the hierarchy are strict, as the supremum norm and the discrete
metric show. Boundedness of the budget set requires strictly positive prices.
Completeness, defined in \autoref{ch:topology}, upgrades normed and inner product
spaces to Banach and Hilbert spaces, the two settings in which
Chapters~\ref{ch:projection} and~\ref{ch:dynamic} respectively work.

\begin{exercise}\label{ex:met-1}
	Prove the polarisation identity
	$\ip{x}{y}=\tfrac14\bigl(\norm{x+y}^2-\norm{x-y}^2\bigr)$ in any inner product
	space, and deduce that an inner product is determined by its norm. Combine this
	with \autoref{rem:parallelogram} to give a complete criterion for when a norm
	comes from an inner product.
\end{exercise}

\begin{exercise}\label{ex:met-2}
	Show that on $\R^n$ the norms $\norm{\cdot}_1$, $\norm{\cdot}_2$ and
	$\norm{\cdot}_\infty$ satisfy
	$\norm{x}_\infty\leq\norm{x}_2\leq\norm{x}_1\leq n\norm{x}_\infty$, so that all
	three induce the same bounded sets and, by \autoref{ch:topology}, the same open
	sets. Explain why the analogous statement fails on $\Cspace{[0,1]}{\R}$ for
	$\norm{f}_1=\int_0^1\abs{f}$ versus $\norm{f}_\infty$, and give a sequence
	witnessing the failure.
\end{exercise}

\begin{exercise}\label{ex:met-3}
	Let $W\in M_{n\times n}(\R)$ be symmetric. Show that
	$\ip{x}{y}_W=x^{\transp}Wy$ is an inner product on $\R^n$ if and only if
	$W\succ0$, and that positive \emph{semi}definiteness gives only a seminorm.
	Identify the subspace on which the seminorm vanishes in terms of $\Ker W$.
\end{exercise}
